\title{The Era of 1-bit LLMs: All Large Language Models are in 1.58 Bits}

\begin{abstract}
The emergence of methods like BitNet~\cite{bitnet} signals the dawn of a new generation of 1-bit Large Language Models (LLMs). This study presents \textbf{\text{BitNet b1.58}}, a novel 1-bit LLM configuration where every parameter (or weight) is restricted to the ternary set \{-1, 0, 1\}. This design achieves parity with conventional full-precision Transformer LLMs (i.e., those using FP16 or BF16) in terms of perplexity and downstream task performance when evaluated with the same model size and quantity of training data. Remarkably, it does so while greatly improving efficiency regarding latency, memory usage, throughput, and energy demands. Importantly, the 1.58-bit LLM establishes a new scaling paradigm and outlines a training methodology for next-generation LLMs that excel in both effectiveness and computational cost. Additionally, it introduces fresh opportunities for developing new computational frameworks and facilitates the creation of dedicated hardware tailored for 1-bit LLM architectures.
\end{abstract}

\section{The Era of 1-bit LLMs}

The advancement of artificial intelligence has been marked by increasingly large and powerful Large Language Models (LLMs), delivering state-of-the-art outcomes across various natural language processing benchmarks. Nevertheless, the rapid expansion in their scale has created obstacles for widespread adoption, particularly due to mounting operational costs and significant environmental considerations arising from their substantial energy demands.

To mitigate these issues, one promising direction has been to apply post-training quantization strategies, yielding lower-bitwidth models primarily for inference~\cite{smoothquant, gptq, quip, quip_sharp}. By constraining the numerical precision of both weights and activations, these techniques dramatically diminish the computational and memory footprint of LLMs. There has been a notable movement from 16-bit representations toward more compact formats, including 4-bit implementations~\cite{gptq, awq}. While prevalent in industry-level LLM deployments, post-training quantization remains inherently sub-optimal, prompting ongoing research into alternative and more effective compression methods.

Advancements in 1-bit model architectures, such as BitNet~\cite{bitnet}, signal a compelling trajectory for diminishing LLM computational costs while largely preserving accuracy. Conventional LLMs operate with 16-bit floating-point representations (e.g., FP16 or BF16), and their core computational load is dominated by matrix multiplications, relying extensively on floating-point additions and multiplications. In contrast, BitNet streamlines this process by restricting matrix multiplications to integer additions alone, thus dramatically curbing the associated energy consumption. Given that power limitations frequently cap the performance of modern chips, these energy gains can also result in increased computational speed.

Beyond raw computation, transferring parameters from off-chip DRAM to the on-chip accelerator’s memory (such as SRAM) during inference constitutes a substantial overhead. While expanding SRAM has been proposed as a means to bolster throughput, it comes with much steeper costs relative to DRAM. 1-bit LLMs are advantageous in this context: they drastically cut down memory usage and reduce bandwidth requirements compared to their full-precision counterparts, yielding improvements in both inference speed and efficiency due to reduced weight-loading overhead from DRAM.

Within this study, we present a novel 1-bit LLM variant, named \textbf{\text{BitNet b1.58}}, in which each parameter adopts one of three discrete values: \{-1, 0, 1\}. Introducing the additional “0” state to the original 1-bit BitNet increases the effective representation to 1.58 bits per parameter in binary terms. \text{BitNet b1.58} upholds all the advantageous attributes of the prior 1-bit BitNet, maintaining its multiplication-free computation framework for matrix multiplications and high potential for hardware-level optimizations. It mirrors the original in terms of energy usage and offers pronounced efficiency gains in memory utilization, computational throughput, and latency relative to FP16-based LLMs. Furthermore, \text{BitNet b1.58} delivers two key improvements: (1) It exhibits enhanced modeling capacity by leveraging feature selection, courtesy of the explicit zero in its weight space, thereby boosting 1-bit LLMs’ overall performance. (2) Our empirical evaluations indicate that, at and above the 3B parameter scale and under consistent configuration conditions (such as model size and training corpus), \text{BitNet b1.58} is able to match the perplexity and downstream task results achieved by FP16 baselines.

\section{BitNet b1.58}

\text{BitNet b1.58} builds upon the BitNet model, a Transformer variant where \emph{nn.Linear} layers are substituted with \emph{BitLinear}. This model is trained entirely from the beginning, employing weights quantized to 1.58 bits and activations quantized to 8 bits. This version introduces several changes compared to the original BitNet, which we detail as follows.

\paragraph{Quantization Function.} We utilize an \emph{absmean} quantization scheme to restrict weight values to the set $\{-1, 0, +1\}$. The procedure involves dividing the weight matrix by its mean absolute value, then mapping each entry to the nearest integer among the permissible values:

\begin{equation}
    \widetilde{W} = \mathrm{RoundClip}(\frac{W}{\gamma + \epsilon}, -1, 1),
\end{equation}

\begin{equation}
    \mathrm{RoundClip}(x, a, b) = \max(a, \min(b, \mathrm{round}(x))),
\end{equation}

\begin{equation}
    \gamma = \frac{1}{nm}\sum_{ij} |W_{ij}|.
\end{equation}

The activation quantization mechanism matches that used in BitNet, with the exception that activations are not pre-scaled to $[0, Q_b]$ before applying non-linearities. Instead, all activations are uniformly rescaled to $[-Q_b, Q_b]$ for each token, which eliminates the requirement for zero-point quantization. This adjustment simplifies both practical implementation and optimization at the system level, while our results indicate that it incurs only a minor impact on performance.

\paragraph{LLaMA-alike Components.}

Given that the LLaMA~\cite{llama, llama2} framework has become a widely adopted foundation for open-source LLMs, our implementation of \text{BitNet b1.58} incorporates similar architectural elements. Namely, it features RMSNorm~\cite{rmsnorm}, SwiGLU~\cite{swiglu}, rotary embedding~\cite{rope}, and the exclusion of bias terms. Consequently, \text{BitNet b1.58} is readily compatible with established open-source platforms, including Huggingface, vLLM~\cite{vllm}, and llama.cpp\footnote{\url{https://github.com/ggerganov/llama.cpp}}, facilitating straightforward integration.

\section{Results}

We conducted a comparison between \text{BitNet b1.58} and our FP16 implementation of \text{LLaMA LLM} across multiple model scales. To provide an equitable assessment, all models were pre-trained using the RedPajama dataset~\cite{redpajama} over a corpus of 100 billion tokens. For evaluation, we assessed zero-shot performance on a suite of language benchmarks including ARC-Easy~\cite{arc}, ARC-Challenge~\cite{arc}, Hellaswag~\cite{hellaswag}, Winogrande~\cite{winoGrande}, PIQA~\cite{piqa}, OpenbookQA~\cite{openbookqa}, and BoolQ~\cite{boolq}. Additionally, validation perplexities were computed on WikiText2~\cite{wikitext2} and C4~\cite{c4} datasets.

We also measured the GPU memory usage and inference latency for both \text{LLaMA LLM} and \text{BitNet b1.58}. These metrics were obtained using the FasterTransformer\footnote{\url{https://github.com/NVIDIA/FasterTransformer}} framework, which offers efficient large language model inference on GPUs. The 2-bit Ladder kernel~\cite{ladder} was employed in the evaluation of \text{BitNet b1.58}. Our main reported metric for latency was the time taken per output token, as this directly impacts inference cost.

As shown in Table~\ref{tab:ppl}, both the perplexity outcomes and resource demands for \text{BitNet b1.58} and \text{LLaMA LLM} are summarized. Notably, at the 3B model size, \text{BitNet b1.58} achieves perplexity comparable to its full-precision counterpart, while offering a 2.71-fold increase in inference speed and reducing GPU memory consumption by a factor of 3.55. For instance, the 3.9B version of \text{BitNet b1.58} outpaces the \text{LLaMA LLM} 3B in speed (by 2.4x), uses significantly less memory (3.32x lower), and delivers stronger performance.

Table~\ref{tab:zero-shot} details the zero-shot accuracy achieved on each downstream task, following the protocol implemented by \emph{lm-evaluation-harness}\footnote{\url{https://github.com/EleutherAI/lm-evaluation-harness}}. The data indicate that as model scale grows, the difference in performance between \text{BitNet b1.58} and \text{LLaMA LLM} diminishes. Crucially, \text{BitNet b1.58} is able to match the zero-shot results of the full precision baseline from the 3B scale onward. This pattern, consistent with the perplexity findings, shows that the 3.9B \text{BitNet b1.58} model surpasses the 3B \text{LLaMA LLM} in both accuracy and efficiency. Collectively, these results position \text{BitNet b1.58} as a Pareto improvement over the prevailing large language models.

\paragraph{Memory and Latency}

We extended our experiments to larger model configurations of 7B, 13B, and 70B parameters, systematically examining the computational costs. As depicted in Figure~\ref{fig:latency-memory-energy}, both latency and memory usage exhibit a scaling trend, where greater model sizes lead to larger speed-ups. Notably, \text{BitNet b1.58} 70B demonstrates a 4.1-fold speed improvement over the \text{LLaMA LLM} baseline. This acceleration is attributed to the increasing time required for the \emph{nn.Linear} operation as the model size grows. Memory consumption mirrors this pattern, given that the embeddings, maintained at full precision, constitute a smaller fraction of total memory usage in larger models. All metrics regarding latency and memory were captured using a 2-bit kernel, implying that further optimizations remain feasible to reduce resource demands.

\paragraph{Energy}

We further conducted energy consumption estimates for arithmetic operations within both \text{BitNet b1.58} and \text{LLaMA LLM}, emphasizing matrix multiplication since it accounts for the primary computational expense in large language models. The breakdown shown in Figure~\ref{fig:energy} indicates that \text{BitNet b1.58} relies predominantly on INT8 addition, in contrast to \text{LLaMA LLM}, which uses both FP16 addition and FP16 multiplication. Referencing the energy model from~\cite{energycost, pokebnn}, we find that \text{BitNet b1.58} reduces the arithmetic operation energy required for matrix multiplication by a factor of 71.4 when implemented on 7nm process technology. Additionally, our measurements of total energy expenditure with input sequences of 512 tokens reveal that the efficiency of \text{BitNet b1.58} relative to the FP16 \text{LLaMA LLM} baseline increases as model size grows. This effect results from the rising contribution of \emph{nn.Linear} in larger models, whereas the impact of other components becomes less significant at scale.

\paragraph{Throughput}

We evaluate and contrast the throughput performance of \text{BitNet b1.58} and \text{LLaMA LLM}, both with 70B parameters, across two 80GB A100 GPUs. Pipeline parallelism~\cite{gpipe} is employed to facilitate the deployment of the \text{LLaMA LLM} 70B model. The batch size was incrementally increased up to the point where GPU memory constraints became a bottleneck, with sequence lengths fixed at 512. As summarized in Table~\ref{tab:throughput}, the 70B parameter \text{BitNet b1.58} model is able to accommodate up to eleven times the batch size relative to \text{LLaMA LLM}, leading to a throughput improvement by a factor of 8.9.

\textbf{\text{BitNet b1.58} enables a revised scaling law connecting inference cost and model effectiveness}. As further context, using insights from Figure~\ref{fig:latency-memory-energy} and \ref{fig:energy}, equivalences between 1.58-bit and 16-bit model capacities can be outlined as follows:

\begin{itemize}
    \item The 13B parameter BitNet b1.58 demonstrates greater efficiency—across latency, memory footprint, and energy consumption—than a 3B parameter FP16 LLM.
    \item BitNet b1.58 with 30B parameters surpasses a 7B FP16 LLM regarding latency, memory utilization, and energy costs.
    \item The 70B BitNet b1.58 model is more efficient by all aforementioned metrics than a 13B FP16 LLM.
\end{itemize}

\paragraph{Training with 2T Tokens}

Token count during pretraining plays a vital role in scaling LLMs. To evaluate how \text{BitNet b1.58} handles a large token budget, we pre-trained a \text{BitNet b1.58} model with 2T tokens, adhering to the data protocol utilized for StableLM-3B~\cite{StableLM-3B-4E1T}, which is currently a leading open-source 3B parameter model. The performance of both models was assessed on a composite benchmark containing Winogrande~\cite{winoGrande}, PIQA~\cite{piqa}, SciQ~\cite{sciq}, LAMBADA~\cite{lambada}, and ARC-easy~\cite{arc}. Zero-shot accuracy results are detailed in Table~\ref{tab:2t}. For tasks reported with both accuracy and normalized accuracy metrics, we computed the mean. Accuracy figures for StableLM 3b at 2T tokens were sourced directly from its original technical documentation. The results indicate that \text{BitNet b1.58} consistently outperforms across all assessed tasks, which highlights robust generalization abilities in 1.58-bit LLMs.

\section{Discussion and Future Work}

\textbf{1-bit Mixture-of-Experts (MoE) LLMs}

Mixture-of-Experts (MoE) architectures have established themselves as an efficient method for scaling LLMs, notably decreasing computation-related FLOPs. However, MoE models continue to grapple with substantial memory usage and substantial inter-chip communication costs, which hinder broader deployment. Transitioning to 1.58-bit LLMs offers viable solutions to these bottlenecks. A lower memory demand reduces the overall device count necessary for MoE deployment. Additionally, network transfer costs for activation data are markedly diminished. Ideally, the memory and communication challenges disappear altogether if the MoE can reside fully within the memory capacity of a single chip.

\textbf{Native Support of Long Sequence in LLMs}

Accommodating long input sequences is increasingly vital for modern LLMs. A prominent barrier to efficient long-sequence inference lies in the heightened memory demands for storing KV caches. \text{BitNet b1.58} takes a promising stride forward by shrinking activation precision from 16 bits to 8 bits, thus allowing a doubling of maximum sequence length for a fixed resource budget. Future investigations could focus on losslessly compressing these representations even further—down to 4 bits or below—in the context of 1.58-bit LLMs, an avenue we intend to pursue.

\textbf{LLMs on Edge and Mobile}

Employing 1.58-bit LLMs presents a significant opportunity for enhancing LLM utility on edge and mobile platforms, which frequently face strict constraints on memory and compute resources. The minimized memory and power requirements of these models facilitate deployment in environments where large-scale models were previously infeasible, broadening the possible use cases. Such improvements can considerably advance the functionality of edge and mobile hardware and drive novel LLM applications. Furthermore, 1.58-bit LLMs are particularly compatible with CPUs, the predominant processors in these settings, meaning that \text{BitNet b1.58} achieves efficient runtime performance, further benefiting device capabilities.

\textbf{New Hardware for 1-bit LLMs}

Recent developments, such as Groq\footnote{\url{https://groq.com/}}, highlight both the feasibility and advantages of developing tailored hardware accelerators (e.g., LPUs) designed to support LLM workloads. Looking forward, we advocate for research into systems and hardware architectures purpose-built for 1-bit LLMs, leveraging the computational efficiencies made possible by the BitNet paradigm~\cite{bitnet}.